% ID: SEL-D0023-OB-E03
% ESTADO: cotejada. Segunda pasada visual completa y borrador compilado revisados.
% Fuente: Selectividad/matematicas_aplicadas/reserva_b_EXAMEN 2 -ANDALUCÍA.pdf
% SHA-256: 8cdb86259405652622f2ea129bffbdb239076ac612fd7fcebb86b9d5bdc451ec
% Página: 2; opción B; ejercicio original 3.

En un aeropuerto internacional operaron 300000 vuelos en un determinado año, distribuidos de la siguiente forma: 150000 en la terminal A, 100000 en la B y 50000 en la C. En ese año se sabe que sufrieron retrasos el 10\% de los vuelos de la terminal A, el 8\% de la B y el 5\% de la C. Determine, para un vuelo elegido al azar, las probabilidades de los siguientes sucesos:
\begin{enumerate}[label=\alph*)]
\item \textbf{(1.25 puntos)} Que no sufriera retraso.
\item \textbf{(1.25 puntos)} Que operase en la terminal A, sabiendo que tuvo retraso.
\end{enumerate}
